\documentclass[11pt]{article}
\usepackage[a4paper,margin=26mm]{geometry}
\usepackage[T1]{fontenc}
\usepackage{lmodern,amsmath,amssymb,amsthm,mathtools,booktabs}
\usepackage[colorlinks=true,linkcolor=blue,citecolor=blue,urlcolor=blue]{hyperref}
\newtheorem{theorem}{Theorem}
\newtheorem{lemma}[theorem]{Lemma}
\newtheorem{corollary}[theorem]{Corollary}
\newtheorem{proposition}[theorem]{Proposition}
\newcommand{\Z}{\mathbb Z}
\newcommand{\Q}{\mathbb Q}
\newcommand{\Rl}{\mathbb Z_{(\ell)}}
\newcommand{\supp}{\operatorname{supp}}
\newcommand{\rad}{\operatorname{rad}}
\newcommand{\OO}{\mathcal O}
\title{The Exact Congruence Content of the Okada Lattice\\\large Simultaneous classification for every residue modulus}
\author{Research note}
\date{15 November 2026}
\begin{document}
\maketitle
\begin{abstract}
For $q>1$, let $\OO_q^{\Z}$ be the lattice of integer-valued, $q$-periodic functions $f$ with $f(0)=0$, zero mean, and $\sum_{n\ge1}f(n)/n=0$. We prove that its image contains the constant-one residue class modulo $M\ge2$ if and only if $M$ divides
\[
 D(q)=\gcd\bigl(q-1,\{v_p(q):p\mid q\}\bigr).
\]
The proof upgrades a first-level prime-adic support elimination to every prime power and then uses an integral Chinese remainder argument. It is constructive in the admissible cases. The result gives nonvanishing, and hence transcendence, for all convergent periodic sums in the excluded congruence classes. The first-level classification already appears in VMRC [056]; the present claim is the higher-power and simultaneous content theorem. Historical priority has not been established. The theorem does not settle the remaining binary sign problem at square periods.
\end{abstract}

\section{The content theorem}
For a rational-valued $q$-periodic function with $f(0)=0$, the series
\[
 L(1,f)=\sum_{n\ge1}\frac{f(n)}n
\]
converges exactly when $\sum_{a=1}^{q-1}f(a)=0$. Let $\OO_q$ denote the rational vector space of such functions for which $L(1,f)=0$, and let $\OO_q^{\Z}=\OO_q\cap\Z^{q-1}$. Write $q=\prod_{p\mid q}p^{e_p}$ and define
\[
 D(q)=\gcd(q-1,e_p:p\mid q).
\]

\begin{theorem}\label{thm:main}
For $q>1$ and $M\ge2$, the following are equivalent:
\begin{enumerate}
\item There is $f\in\OO_q^{\Z}$ with $f(a)\equiv1\pmod M$ for every $1\le a<q$.
\item $q\equiv1\pmod M$ and $M\mid e_p$ for every prime $p\mid q$.
\item $M\mid D(q)$.
\end{enumerate}
Equivalently, the content of the class of $\boldsymbol 1$ in the free abelian group $\Z^{q-1}/\OO_q^{\Z}$ is exactly $D(q)$.
\end{theorem}

Here the content of an element of a free abelian group is the largest positive integer by which that element is divisible. The quotient is torsion-free because $\OO_q^{\Z}$ is the integer kernel of rational linear equations. Its constant-one class is nonzero, since the constant-one vector has nonzero mean. Divisibility of this class by $M$ is exactly $\boldsymbol 1\in\OO_q^{\Z}+M\Z^{q-1}$.

The prime-power version is the central assertion.
\begin{theorem}\label{thm:local}
Let $\ell$ be prime and $k\ge1$. There exists $f\in\OO_q$ with
\[
 f(a)\in1+\ell^k\Rl\qquad(1\le a<q)
\]
if and only if
\begin{equation}\label{eq:localconditions}
 q\equiv1\pmod{\ell^k},\qquad \ell^k\mid e_p\quad(p\mid q).
\end{equation}
These conditions are also equivalent to $q=m^{\ell^k}$ for an integer $m\equiv1\pmod\ell$.
\end{theorem}

The last equivalence follows by reduction modulo $\ell$ in one direction and the elementary valuation bound
\[
 v_\ell(x^a-1)\ge v_\ell(x-1)+v_\ell(a)\qquad(x\equiv1\pmod\ell)
\]
in the other. This lower bound also holds for $\ell=2$ and odd $x$.

\section{Okada's equations and the prime-adic quotient}
We recall the rational vanishing criterion used in the proof; these equations are classical \cite{Okada,MP,Bharadwaj}. For $d\mid q$, $d>1$, put $U_d=(\Z/d\Z)^\times$ and
\[
 f_d(x)=f((q/d)x),\qquad S_d=\sum_{x\in U_d}f_d(x).
\]
Together with mean zero, the criterion consists of the prime equations
\begin{equation}\label{eq:prime}
 \sum_{d\mid q,\ d>1}\left(v_p(q/d)-\frac{[p\mid d]}{p-1}\right)S_d=0
 \qquad(p\mid q)
\end{equation}
and, for each nonprincipal primitive Dirichlet character $\chi$ of conductor $m\mid q$, the character equation
\begin{equation}\label{eq:character}
 \sum_{m\mid d\mid q}\frac d{\varphi(d)}
 \prod_{b\mid d,\ b\nmid m}\left(1-\frac{\chi(b)}b\right)
 \sum_{x\in U_d}f_d(x)\overline{\chi(x)}=0.
\end{equation}
Equivalently one may replace the character equations by
\begin{equation}\label{eq:unit}
 \sum_{a\in\mathcal M(q)}\frac{f(ua)}a=0\qquad(u\in U_q),
\end{equation}
where $\mathcal M(q)$ is the monoid generated by the prime divisors of $q$.

Mean zero transforms \eqref{eq:prime} into
\begin{equation}\label{eq:Bstar}
 B_p^*:=\sum_{d:p\mid d}\bigl((p-1)v_p(d)+1\bigr)S_d=0.
\end{equation}
We will repeatedly use the elementary identities
\begin{align}
 \sum_{d:p\mid d}\varphi(d)&=\frac q{p^{e_p}}(p^{e_p}-1),\label{eq:divsum1}\\
 \sum_{d:p\mid d}v_p(d)\varphi(d)
 &=\frac q{p^{e_p}}\left(e_pp^{e_p}-\frac{p^{e_p}-1}{p-1}\right),\label{eq:divsum2}\\
 \sum_{d:p\mid d}\bigl((p-1)v_p(d)+1\bigr)\varphi(d)
 &=e_p(p-1)q.\label{eq:divsum3}
\end{align}
They follow by separating the $p$-part of a divisor and using $\sum_{d\mid n}\varphi(d)=n$.

Fix $\ell$, and suppose that $f\in\OO_q$ is $1$ modulo $\ell^k$. Mean zero implies $q\equiv1\pmod{\ell^k}$, so $\ell\nmid q$. The Okada space is invariant under multiplication of the argument by units: this scales each character equation by a root of unity and preserves the prime equations. Average $f$ over the subgroup of $U_q$ of order prime to $\ell$. Its order is a unit of $\Rl$, so the resulting function is still $1$ modulo $\ell^k$.

For each $p\mid q$, set
\[
 N_p=\ell^{v_\ell(p-1)},\qquad Q_p=C_{N_p},\qquad
 P_+=\{p:N_p>1\},\quad P_0=\{p:N_p=1\}.
\]
The maximal $\ell$-quotient of $U_d$ is $Q_d=\prod_{p\mid d}Q_p$. Its fibres have size $h_d=\varphi(d)/\prod_{p\mid d}N_p$, a unit of $\Rl$. Let $\Phi_d$ be the actual sum of $f_d$ over each fibre. Then
\begin{equation}\label{eq:fiberres}
 \Phi_d(y)\equiv h_d\pmod{\ell^k\Rl},\qquad
 S_d\equiv\varphi(d)\pmod{\ell^k\Rl}.
\end{equation}
The first congruence concerns actual sums, not normalized averages; this distinction is essential at higher powers.

For every nonempty set $S$ of prime divisors of $q$, define
\[
 G_S=\sum_{\supp(d)=S}\Phi_d\quad\text{on }Q_S,
 \qquad H_p=\frac{p^{e_p}-1}{N_p}.
\]
Summing the fibre sizes prime by prime gives the exact residue formula
\begin{equation}\label{eq:GSres}
 G_S\equiv\prod_{p\in S}H_p\pmod{\ell^k\Rl}.
\end{equation}
For a character nontrivial on each factor of $Q_A$, its primitive support is $A\subseteq P_+$. The normalized retained character equations are
\begin{equation}\label{eq:compressedA}
 \sum_{S\supseteq A}\prod_{b\in S\setminus A}\frac{b-\chi(b)}{b-1}
 \widehat G_S(\chi)=0.
\end{equation}

\section{Folding and integral support elimination}
Let $\tau_b$ denote translation by the image of $b$ on the relevant quotient; the convention is chosen so that $\widehat{\tau_bF}(\chi)=\chi(b)\widehat F(\chi)$. For nonempty $S\subseteq P_+$ put
\begin{equation}\label{eq:fold}
 \widetilde G_S=\sum_{U\subseteq P_0}
 \left(\prod_{b\in U}\frac{b-\tau_b}{b-1}\right)G_{S\cup U},
 \qquad r=\prod_{b\in P_0}b^{e_b}.
\end{equation}
Every folding operator is $\Rl$-integral, preserves total sums, and acts as the identity on a constant residue modulo $\ell^k$. Thus \eqref{eq:compressedA} becomes the same system using only primes in $P_+$, and
\begin{equation}\label{eq:foldres}
 \widetilde G_S\equiv r\prod_{p\in S}H_p\pmod{\ell^k\Rl}.
\end{equation}
For $p\in P_+$ define $T_p(G)=\sum_{S\ni p}\sum_{y\in Q_S}G_S(y)$. Folding leaves this quantity unchanged. Before folding, \eqref{eq:Bstar} gives
\begin{equation}\label{eq:Toriginal}
 T_p(G)=-(p-1)\sum_{d:p\mid d}v_p(d)S_d.
\end{equation}
In particular, $T_p(G)/N_p$ is $\Rl$-integral.

We give the support-elimination details so the higher-power argument has no implicit integrality assumption.
\begin{lemma}[Anchored decomposition]\label{lem:anchor}
Let $|S|\ge2$ and let $F:Q_S\to\Rl$ have no Fourier component nontrivial in every coordinate. Then
\[
 F=\sum_{p\in S}\operatorname{inf}_{S\setminus\{p\}}^S h_p
\]
with $\Rl$-valued functions $h_p$. If $F$ has constant residue modulo an integer power of $\ell$, all $h_p$ can be chosen with constant residue. If $F$ is zero modulo that power, all $h_p$ can be chosen zero modulo that power.
\end{lemma}
\begin{proof}
Use anchored inclusion--exclusion at the identity in each coordinate. The top interaction is zero because the Fourier expansion has no term nontrivial in all coordinates. Each remaining interaction omits a coordinate and can be assigned to one such omitted coordinate. Every nonempty anchored difference of a constant-residue function is zero modulo the indicated power. The empty interaction is constant. No division is involved.
\end{proof}

\begin{lemma}[Elimination circuit]\label{lem:circuit}
For $S\subseteq P_+$, $|S|\ge2$, $p\in S$, and $U=S\setminus\{p\}$, the changes
\[
 \delta G_S=-\operatorname{inf}_U^S h,
 \qquad \delta G_U=\frac{N_p}{p-1}(p-\tau_p)h
\]
preserve all equations \eqref{eq:compressedA}. They preserve $T_b$ for $b\ne p$ and change $T_p$ by $-N_p\sum_{Q_U}h$.
If $h\equiv0\pmod{\ell^k}$, they preserve every lower-layer residue modulo $\ell^k$ and every $T_b$ modulo $\ell^kN_b$. If $k=1$ and $h$ merely has constant residue, the same preservation statements hold modulo $\ell$ and $\ell N_b$.
\end{lemma}
\begin{proof}
The coefficient $N_p/(p-1)$ is a unit of $\Rl$. For characters whose support contains $p$, inflation has zero Fourier coefficient. For support contained in $U$, inflation contributes a factor $N_p$, and the two altered terms cancel because their multipliers are $-N_p(p-\chi(p))/(p-1)$ and its opposite. Total sums in the two layers change by $-N_p\sum h$ and $+N_p\sum h$, proving the assertions about $T_b$.

The higher-power claims are immediate if $h$ is zero modulo $\ell^k$. At the first level, $p\equiv1\pmod\ell$ makes $(p-\tau_p)h$ zero modulo $\ell$ for constant-residue $h$. Also $|Q_U|$ is divisible by $\ell$, so $\sum h\in\ell\Rl$.
\end{proof}

Apply these circuits in decreasing support size. When a support $S$ is reached, all its strict supersets have vanished. Its character equation says exactly that its fully nontrivial Fourier component is zero. Lemma~\ref{lem:anchor} and the circuits eliminate it, altering only smaller supports. At termination each singleton is a constant function $a_p$.

\section{Necessity at every prime power}
We prove Theorem~\ref{thm:local} by induction on $k$. The condition $q\equiv1\pmod{\ell^k}$ already follows from the mean. For any $p\in P_0$, equations \eqref{eq:Bstar}, \eqref{eq:fiberres}, and \eqref{eq:divsum3} give
\[
 0\equiv e_p(p-1)q\pmod{\ell^k}.
\]
Both $(p-1)$ and $q$ are units, so $\ell^k\mid e_p$ without induction.

\subsection{First level}
For $p\in P_+$, every $S_d$ with $p\mid d$ is zero modulo $\ell$. Equation~\eqref{eq:Toriginal} implies
\[
 T_p/N_p\equiv0\pmod\ell.
\]
The folded functions have constant residues by \eqref{eq:foldres}. First-level elimination preserves the singleton residues and the preceding congruences. Since
\[
 H_p\equiv e_p\frac{p-1}{N_p}\pmod\ell,
\]
the terminal constant satisfies
\[
 0\equiv a_p\equiv r e_p\frac{p-1}{N_p}\pmod\ell.
\]
The two factors besides $e_p$ are units, so $\ell\mid e_p$. This reproduces the first-level result of VMRC [056] \cite{VMRC}, rather than treating it as a new conclusion.

\subsection{Higher levels}
Suppose $k\ge2$. The induction hypothesis applies to the same function modulo $\ell^{k-1}$ and gives $\ell^{k-1}\mid e_p$ for all $p$. For $p\in P_+$, the valuation bound for $p^{e_p}-1$ implies
\[
 H_p\in\ell^{k-1}\Rl.
\]
Hence \eqref{eq:foldres} is zero modulo $\ell^k$ for every support of size at least two, because $2(k-1)\ge k$. Eliminate these supports using functions $h$ that are all zero modulo $\ell^k$. The terminal singleton constants therefore satisfy
\begin{equation}\label{eq:terminal}
 a_p\equiv rH_p\pmod{\ell^k},\qquad
 a_p\equiv T_p(G)/N_p\pmod{\ell^k}.
\end{equation}
The second congruence uses preservation modulo $\ell^kN_p$ before division by $N_p$.

Fix $p\in P_+$, abbreviate $e=e_p$, and put
\[
 u_p=\frac{p-1}{N_p},\qquad A_e(p)=\frac{p^e-1}{p-1}=\sum_{j=0}^{e-1}p^j.
\]
Equation~\eqref{eq:Toriginal}, divided by $N_p$ before reduction, and \eqref{eq:divsum2} give
\[
 \frac{T_p(G)}{N_p}
 \equiv-u_p\frac q{p^e}\bigl(ep^e-A_e(p)\bigr)\pmod{\ell^k}.
\]
Combining this with \eqref{eq:terminal} and $H_p=u_pA_e(p)$, then cancelling the unit $u_p$, yields
\begin{equation}\label{eq:key}
 e q\equiv\left(\frac q{p^e}-r\right)A_e(p)\pmod{\ell^k}.
\end{equation}
For each other $a\in P_+$, $a\equiv1\pmod\ell$ and $\ell^{k-1}\mid e_a$, so $a^{e_a}\equiv1\pmod{\ell^k}$. Therefore
\[
 \frac q{p^e}-r=r\left(\prod_{a\in P_+\setminus\{p\}}a^{e_a}-1\right)
 \equiv0\pmod{\ell^k}.
\]
Equation~\eqref{eq:key} and invertibility of $q$ imply $\ell^k\mid e$. This completes necessity for every prime divisor.

\section{Explicit constructions for sufficiency}
Assume \eqref{eq:localconditions}, and write $m=\ell^k$.

\subsection{A prime power with $p\equiv1\pmod\ell$}
Let $m\mid e$. The integer $A_e(p)$ is divisible by $m$, as is
\[
 C=A_e(p)-e=\sum_{j=1}^{e-1}j\varphi(p^{e-j}).
\]
Define $g$ first on the nonzero multiples of $p$ modulo $p^e$: set all values to $1$, except $g(p)=1-C$. On each unit $a$, define
\begin{equation}\label{eq:primeconstruct}
 g(a)=-\sum_{j=1}^{e-1}\frac{g(p^ja)}{p^j},\qquad g(0)=0.
\end{equation}
All denominators are prime to $\ell$. The nonunit values are $1$ modulo $m$, and
\[
 g(a)\equiv-\sum_{j=1}^{e-1}p^{-j}\equiv1\pmod m,
\]
because $\sum_{j=0}^{e-1}p^{-j}=p^{1-e}A_e(p)\equiv0\pmod m$.

The unit equations \eqref{eq:unit} hold by construction. Let $V_j$ be the sum of $g$ over residues of valuation $j$. The map from units to the valuation-$j$ layer has $p^j$ preimages for each residue, so \eqref{eq:primeconstruct} gives $V_0=-\sum_{j\ge1}V_j$, proving mean zero. Moreover,
\[
 \sum_{j=1}^{e-1}jV_j=C-C=0.
\]
Under mean zero this is exactly the prime equation. Thus $g\in\OO_{p^e}$.

\subsection{The primes with $\ell\nmid p-1$}
Set $r=\prod_{b\in P_0}b^{e_b}$. If $r=1$, omit this factor. Otherwise $r\equiv1\pmod m$, since $q\equiv1\pmod m$ and every prime-power factor belonging to $P_+$ is $1$ modulo $m$. For every divisor $d>1$ of $r$, $\varphi(d)$ is a unit of $\Rl$.

Construct a function constant on each gcd layer, beginning with layer totals $\Phi_d=\varphi(d)$. Their mean is $r-1$, and the prime totals in \eqref{eq:Bstar} are $e_b(b-1)r$. Choose $b_0\mid r$ and subtract $r-1$ from $\Phi_{b_0}$. The mean becomes zero, and the new prime totals are
\[
 \beta_b=e_b(b-1)r-[b=b_0]b(r-1)\in m\Rl.
\]
For each $b\mid r$, put $w_b=-\beta_b/(b-1)\in m\Rl$, subtract $w_b$ from $\Phi_b$, and add it to $\Phi_{b^2}$. These layers exist because $e_b\ge m\ge2$. This preserves the mean, changes $B_b^*$ by $(b-1)w_b=-\beta_b$, and leaves the other prime totals unchanged.

Assign the constant value $\Phi_d/\varphi(d)$ to the layer of denominator $d$. These values are all $1$ modulo $m$. Their nonprincipal character coefficients are zero, and the mean and prime equations have just been made zero. This produces an element of $\OO_r$ in the required congruence class.

\subsection{Composition and integral Chinese remainders}
For periods $A,B$ and functions $u\in\OO_A$, $v\in\OO_B$, both $1$ modulo $m$ off zero, define
\[
 F(n)=\begin{cases}u(n),&A\nmid n,\\v(n/A),&A\mid n.\end{cases}
\]
It has period $AB$, zero mean, the same congruence class, and
\[
 L(1,F)=L(1,u)+A^{-1}L(1,v)=0.
\]
Composing the constructed prime-power and $P_0$ factors proves sufficiency in Theorem~\ref{thm:local}.

To pass to integers, clear the denominators of a local solution with an integer $d$ prime to $\ell$, then multiply by an integer $t$ with $td\equiv1\pmod{\ell^k}$. This gives an integer Okada vector in the same local congruence class.

Finally let $M=\prod_i\ell_i^{k_i}$. Necessity in Theorem~\ref{thm:main} follows by applying Theorem~\ref{thm:local} at each prime. Conversely, construct integer vectors $g_i\in\OO_q^{\Z}$ with $g_i\equiv1\pmod{\ell_i^{k_i}}$. Choose integer Chinese-remainder coefficients $c_i$ with $c_i\equiv1$ modulo the $i$th prime power and $c_i\equiv0$ modulo all others. Then
\[
 f=\sum_i c_i g_i\in\OO_q^{\Z},\qquad f\equiv\boldsymbol1\pmod M.
\]
This proves Theorem~\ref{thm:main} in full.

\section{Examples and consequences}
\begin{corollary}[Least admissible period]
For a prime-power residue modulus $M=\ell^k$, the least period $q>1$ admitting an integer vanishing function that is $1$ modulo $M$ off zero is
\[
 q_{\min}=(\ell+1)^{\ell^k}.
\]
\end{corollary}
\begin{proof}
Theorem~\ref{thm:local} requires $q=a^{\ell^k}$ with $a>1$ and $a\equiv1\pmod\ell$. Its least possible base is $\ell+1$, and sufficiency supplies an integer witness there.
\end{proof}
Thus the least periods for residue moduli $2,4,8$ are respectively $9,81,6561$, and those for residue moduli $3,9$ are $64,262144$. These statements concern unrestricted integer coefficients in one residue class, not binary sign functions.

For $q=81$ and $M=4$, the following integer function is an exact witness:
\[
\begin{array}{c|r}
\text{residue class modulo }81& f(a)\\\hline
0&0\\
1,28,55&-311\\
3&945\\
\text{other units}&13\\
\text{other nonzero multiples of }3&-27
\end{array}
\]
All nonzero-residue values are $1$ modulo $4$. The mean is
\[
 3(-311)+51(13)+945+25(-27)=0.
\]
The unit equations have only two forms:
\[
 -311+945/3-27/9-27/27=0,
 \qquad 13-27/3-27/9-27/27=0.
\]
The prime equation reduces to
\[
 -27(18+2\cdot6+3\cdot2)+(945+27)=0.
\]
Thus the witness is certified without numerical evaluation of its infinite sum.

There is also a compact family for $q=2^e$, $e\ge2$. Set all nonzero-residue values to $1$ except
\[
 f(q/2)=(e-3)q+4,\qquad f(q/4)=f(3q/4)=-\frac{(e-2)q}{2}.
\]
It belongs to $\OO_q^{\Z}$ and is $1$ modulo every divisor of $\gcd(q-1,e)$. Its mean is zero. The unit equation is
\[
 \sum_{j=0}^{e-3}2^{-j}
 -2(e-2)+2(e-3)+8/q=0
\]
with the empty-sum interpretation for $e=2$, and the valuation-weighted prime equation is zero by direct expansion. For $q=2^{18}$ this gives $M=9$ and the exceptional values $3932164,-2097152,-2097152$.

\begin{center}
\begin{tabular}{rlll}
\toprule
$q$&factorization&$D(q)$&admissible $M\ge2$\\\midrule
9&$3^2$&2&2\\
64&$2^6$&3&3\\
81&$3^4$&4&2,4\\
225&$3^25^2$&2&2\\
625&$5^4$&4&2,4\\
4096&$2^{12}$&3&3\\
15625&$5^6$&6&2,3,6\\
262144&$2^{18}$&9&3,9\\\bottomrule
\end{tabular}
\end{center}

\begin{corollary}[Uniform nonvanishing and transcendence]
Suppose $M\nmid D(q)$. Every integer-valued, $q$-periodic function $f$ with $f(0)=0$, $f(a)\equiv1\pmod M$ off zero, and zero mean satisfies
\[
 L(1,f)\ne0.
\]
Every such value is transcendental.
\end{corollary}
\begin{proof}
Vanishing would contradict Theorem~\ref{thm:main}. A convergent periodic Dirichlet series with algebraic coefficients at $1$ is an algebraic linear combination of logarithms of algebraic numbers, so Baker's theorem makes a nonzero value transcendental \cite{BBW,Baker}.
\end{proof}
If $M\nmid q-1$, the zero-mean class itself is empty. The substantive nonvanishing cases therefore have $M\mid q-1$ but $M\nmid e_p$ for some prime divisor. For example, at $q=9$, every convergent integer-coefficient sum with all nonzero-residue coefficients $1$ modulo $4$ is transcendental, whereas the displayed period-$81$ example vanishes in that same residue class.

\section{Scope and verification status}
VMRC [056] already gives the $k=1$ theorem for every prime $\ell$, including the odd-square criterion at $\ell=2$. The new step proposed here is preservation at every $\ell^k$, followed by the simultaneous integer content formula. All first-level ingredients have been reproduced above. The higher-power congruences were independently checked, including the factor $h_d$ in the actual orbit sums, the modulus $\ell^kN_p$ before division, and the folding of primes with trivial $\ell$-quotient.

Exact saturated-kernel calculations at $q=9,25,49,81,121,169,225,289,361,441,625,729$ agreed with the predicted maximum $2$-adic depth. These finite checks supplement the proof and are not its basis. The explicit period-$81$ witness and the power-of-two family admit the short full-equation checks displayed above.

The theorem concerns unrestricted integer values in a prescribed congruence class. At $M=2$ the values are odd integers, not necessarily the two signs $\pm1$. It therefore does not settle the remaining square-period sign-lifting problem. It also does not establish irrationality of Euler's constant or a new result about unrelated constants. A complete historical novelty search and external mathematical review remain outstanding.

\begin{thebibliography}{9}
\bibitem{Okada} T. Okada, \emph{On a certain infinite series for a periodic arithmetical function}, Acta Arithmetica \textbf{40} (1982), 143--153.
\bibitem{MP} M. R. Murty and S. Pathak, \emph{The Okada space and vanishing of $L(1,f)$}, Functiones et Approximatio Commentarii Mathematici \textbf{66} (2022), 35--57. DOI:10.7169/facm/1952.
\bibitem{Bharadwaj} A. Bharadwaj, \emph{On primitivity and vanishing of Dirichlet series}, arXiv:2204.03674 (2022). \url{https://arxiv.org/abs/2204.03674}.
\bibitem{BBW} A. Baker, B. J. Birch and E. A. Wirsing, \emph{On a problem of Chowla}, Journal of Number Theory \textbf{5} (1973), 224--236.
\bibitem{Baker} A. Baker, \emph{Transcendental Number Theory}, Cambridge University Press, 1975.
\bibitem{VMRC} QEEDDD, \emph{Verifiable Mathematical Reasoning Corpus}, manuscripts [055]--[057], especially [056], \emph{Square moduli, sparse masks, and sign obstructions for vanishing Erd\H{o}s sums}. \url{https://github.com/QEEDDD/Verifiable-Mathematical-Reasoning-Corpus}.
\end{thebibliography}
\end{document}
